Let $m\ge 1$ be an integer. Let $\mathcal{G}_m$ be the family of finite simple graphs with exactly $m$ edges and no isolated vertices. For $m=\binom{n}{2}+t$ with $n$ maximal and $0\le t < n$, define $H_m$ to be the graph obtained from $K_n$ by adding one new vertex joined to exactly $t$ vertices of the $K_n$ (note: when $t=0$, this construction produces an isolated vertex; since the quantification restricts to graphs with no isolated vertices, the natural comparison graph is $K_n$ itself, but in the counterexample below $t\neq 0$ so no ambiguity arises).

Conjecture (as stated): For every $G\in\mathcal{G}_m$,
\[
R(G)\le R(H_m),
\]
where $R(\cdot)$ denotes the ordinary (two-colour) Ramsey number of a graph.

I will disprove this conjecture by exhibiting $m=2$ and a graph $G\in\mathcal{G}_2$ with $R(G)>R(H_2)$.
